% Generated draft: requires mathematical and source review.
\documentclass[11pt]{article}
\usepackage[margin=25mm]{geometry}
\usepackage{fontspec}
\setmainfont{Latin Modern Roman}
\usepackage{amsmath,amssymb,mathtools,unicode-math}
\setmathfont{Latin Modern Math}
\usepackage{enumitem,needspace,xurl,hyperref}
\setlength{\parindent}{0pt}
\setlength{\parskip}{5pt}
\setlength{\emergencystretch}{3em}
\newcommand{\N}{\mathbb{N}}
\newcommand{\Z}{\mathbb{Z}}
\newcommand{\Q}{\mathbb{Q}}
\newcommand{\R}{\mathbb{R}}
\newcommand{\C}{\mathbb{C}}
\newcommand{\F}{\mathbb{F}}
\newcommand{\A}{\mathbb{A}}
\newcommand{\PP}{\mathbb{P}}
\newcommand{\E}{\mathbb{E}}
\newcommand{\eps}{\varepsilon}
\newcommand{\dd}{\,\mathrm{d}}
\begin{document}
\section*{UnsolvedMath ID: 30004830}
\textbf{Source problem code:} catalogue entry
\subsection{Definitions and mathematical statement}
Work over \(\C\). For a group \(G\), a left-left Yetter--Drinfeld module \(V\in{}^{\C G}_{\C G}\mathcal{YD}\) is a \(G\)-graded vector space \(V=\bigoplus_{g\in G}V_g\) equipped with a linear left \(G\)-action satisfying \(h\cdot V_g\subseteq V_{hgh^{-1}}\). Its support is \(\operatorname{supp}V=\{g\in G:V_g\ne0\}\), and its braiding is
\[
c(v\otimes w)=g\cdot w\otimes v\qquad(v\in V_g).
\]
No finite-dimensionality assumption on \(V\) is imposed in the two principal questions below.

Let \(T(V)=\bigoplus_{n\ge0}V^{\otimes n}\) be the tensor algebra with its braided Hopf algebra structure for which \(V\) is primitive. A pre-Nichols algebra of \(V\) is a quotient \(R=T(V)/J\), where \(J\) is a homogeneous braided Hopf ideal in the Yetter--Drinfeld category and \(J\subseteq\bigoplus_{n\ge2}V^{\otimes n}\). The Nichols algebra is \(\mathcal B(V)=T(V)/I(V)\), where \(I(V)\) is the largest such ideal. Thus there are canonical graded braided Hopf algebra surjections \(T(V)\twoheadrightarrow R\twoheadrightarrow\mathcal B(V)\) that identify degree one with \(V\). Write \(R\#\C G\) for the bosonization; its underlying vector space is \(R\otimes\C G\), and its multiplication satisfies
\[
(r\#g)(s\#h)=r(g\cdot s)\#gh.
\]
For a unital \(\C\)-algebra \(A\), use the convention
\[
\operatorname{GKdim}A=\sup_W\limsup_{n\to\infty}\frac{\log\dim_{\C}W^n}{\log n},
\]
where \(W\) ranges over finite-dimensional vector subspaces of \(A\) containing \(1\), and \(W^n\) is the span of products of \(n\) elements of \(W\).

\textbf{Finite-growth pre-Nichols classification, source Question 3.} For every finitely generated group \(G\) and every \(V\in{}^{\C G}_{\C G}\mathcal{YD}\) satisfying
\[
\operatorname{GKdim}\bigl(\mathcal B(V)\#\C G\bigr)<\infty,
\]
classify all pre-Nichols algebras \(R\) of \(V\) for which
\[
\operatorname{GKdim}\bigl(R\#\C G\bigr)<\infty.
\]
Equivalently, one may seek all defining ideals \(J\) as above satisfying this growth condition. The source does not specify the equivalence relation intended for an isomorphism classification. The finite-generation hypothesis on \(G\) is explicit in the supplied cleaned statement; it is not repeated in the extracted wording of Question 3.

\textbf{Type-C growth assertion, source Conjecture 1, attributed there to reference [2].} Independently of the preceding finite-growth hypothesis, let \(G\) be any group, let \(\mathcal O\) be a finite conjugacy class of \(G\) of type C, and let \(V\in{}^{\C G}_{\C G}\mathcal{YD}\). The proposed assertion is
\[
\operatorname{supp}V=\mathcal O
\quad\Longrightarrow\quad
\operatorname{GKdim}\mathcal B(V)=\infty.
\]
Here the rack operation on a conjugacy class is \(x\triangleright y=xyx^{-1}\). The supplied extract refers to reference [6] for the definition of a finite rack of type C, but provides neither that definition nor the bibliographic identity of [6]. Consequently, ``type C'' remains an unresolved source-defined predicate in this draft; the assertion cannot be made self-contained from the supplied evidence. The original conjecture explicitly allows arbitrary groups. It must not be read as also imposing the finite-GK hypothesis from Question 3.

\textbf{Source-reported context and qualifications.} The extract reports as Theorem 1, attributed to [2] and described as using a theorem of Malcev, that if \(\operatorname{supp}V\) contains an infinite conjugacy class, then \(\operatorname{GKdim}(\mathcal B(V)\#\C G)=\infty\). It then states that every noncentral conjugacy class in a finitely generated torsion-free nilpotent group is infinite, so finite GK dimension of this bosonization forces \(\operatorname{supp}V\subseteq Z(G)\). For finite groups, it reports a result from [6] that support equal to a type-C conjugacy class implies \(\dim_{\C}\mathcal B(V)=\infty\); this is a vector-space dimension statement, distinct from the proposed GK-dimension assertion. The extract credits its proof to results of I. Heckenberger, H.-J. Schneider, and L. Vendramin, based on Weyl groupoid theory referenced as [11]. It also reports, citing [2], that if \(G\) is finitely generated and nilpotent and its nontrivial torsion subgroup has odd order, then every finite conjugacy class is either abelian or of type C. Here an abelian conjugacy class means that its elements commute pairwise. These are supplied contextual claims, not independently verified results of this draft.

The subsequent survey passage separately assumes that \(G\) is abelian and \(V\) is finite-dimensional, and fixes a decomposition \(V=V_1\oplus\cdots\oplus V_\theta\) into indecomposable Yetter--Drinfeld objects. It calls an indecomposable object a point when its dimension is one, a block when its dimension exceeds one and it is indecomposable as a braided vector space, and a pale block when its dimension exceeds one and it is decomposable as a braided vector space. A direct sum of points is called diagonal type. In this setting the extract records historical Conjecture 2, attributed to [3]: \(\operatorname{GKdim}\mathcal B(V)<\infty\) if and only if \(V\) has finite root system, with classification then attributed to [12]. The extract neither defines that root-system condition nor supplies the bibliography for these numbered references, and ends mid-discussion. These later abelian and finite-dimensional hypotheses do not automatically restrict Question 3 or Conjecture 1. The supplied dated literature assessment reports that the diagonal finite-GK conjecture and several low-dimensional or block cases have been settled, while the general pre-Nichols classification and type-C conjecture remain open; this assessment requires independent verification.
\subsection{Short English statement}
For each finitely generated group \(G\) and Yetter--Drinfeld module \(V\) whose Nichols bosonization has finite Gelfand--Kirillov dimension, classify all pre-Nichols algebras whose bosonizations also have finite Gelfand--Kirillov dimension. A separate conjecture concerns arbitrary groups: finite conjugacy-class support of type C should force the Nichols algebra itself to have infinite Gelfand--Kirillov dimension. The supplied record omits the definition of type C and the intended classification equivalence. Its partially solved status and literature assessment are supplied curation, not independently verified conclusions.
\subsection{Sources}
\begin{itemize}
\item Local UnsolvedMath registry, supplied complete source record, permanent ID \(30004830\), title ``Classification of Finite-Growth Pre-Nichols Algebras'', exact source code \texttt{catalogue entry}, collection ``Oberwolfach Reports Open Problems''. No individual proposer is supplied; the proposed year is recorded as 2021.
\item Supplied report citation: \emph{Quantum Groups -- Algebra, Analysis and Category Theory} (2022), Oberwolfach Reports, pp. 2397--2458, DOI \texttt{10.4171/owr/2021/44}. \url{https://doi.org/10.4171/owr/2021/44}. The record describes this collection as workshop reports and open-problem discussions. Its supplied verification note says the report was inspected and Question 3 and the explicit type-C conjecture were isolated from intervening theorem and survey prose; this draft does not independently verify that inspection.
\item Supplied literature source, ``Oberwolfach Report 44'': \url{https://doi.org/10.4171/owr/2021/44}.
\item Supplied literature source, ``Yang on finite-GK pre-Nichols algebras'': \url{https://arxiv.org/abs/2504.12206}.
\item Supplied literature source, ``Angiono on nonabelian pointed Hopf algebras'': \url{https://doi.org/10.1112/plms.12559}.
\item Literature assessment dated exactly \texttt{2026-08-22}: the supplied status is ``partially solved''. The supplied assessment says that the diagonal finite-GK conjecture and several low-dimensional or block cases have been settled, while the type-C conjecture and general finite-GK pre-Nichols classification remain open. This is supplied dated curation derived from the OWR record, not an independently verified literature assessment; the record expressly requests independent verification of time-sensitive status and cited claims.
\item Internal source labels [2], [3], [6], [11], and [12] occur in the extract without their bibliographic entries. In particular, [6] is the supplied pointer for the missing type-C definition. No external bibliographic identities for these labels are inferred here.
\end{itemize}
\end{document}
